% !TeX root = linear-codes.tex
\section{The transition probability}

\subsection{Fourier analysis on the cube}
For functions $f,g \colon\field_2^n\to\reals$, let
$\iprod{f}{g}=\sum_x f(x)g(x)$. Define the Fourier basis
functions $\chi_y$ for $y\in\field_2^n$ as
$\chi_y(x) = (-1)^{x\cdot y}$. Note that
$\iprod{\chi_y}{\chi_y}=2^n$ and for $y\neq z$ we have
$\iprod{\chi_y}{\chi_z}=0$.
We define the Fourier transform of a function
as the coefficients in this character basis:
$$
\hat f(y) = \iprod{f}{\chi_y}.
$$
Note our Fourier transform is $2^{n/2}$ times a unitary.
The inverse transform is thus $2^{-n/2}$ times a unitary
given by
$$
f = \E_y\sparen{\hat f(y)\chi_y}.
$$

\subsection{Fourier estimates for Hamming balls}

The following two consequences of the sharp moment and tail
bounds of Kirshner and Samorodnitsky \cite[Corollary 1.4 and
Theorem 1.5]{KirshnerS2021} are all that we need.  Their
derivation is given in \autoref{sec:moment-tail-derivation}.

\begin{lemma}[Fourier estimates for the punctured Hamming
ball]
\label{lem:ball-fourier-estimates}
Let $w=d-1$, and let
$$
\mu(x)=\frac{\indicate{x\in\ball}}{V}.
$$
Fix $\eps>0$ and suppose
$$
\lg n\ll w\le(1/2-\eps)n.
$$
Then
$$
\E_y\abs{\hat\mu(y)}^3
=O\left(\frac1{Vn^3}\right)
$$
and
$$
\Pr_y\sparen{\abs{\hat\mu(y)}>1/n}
\le
2^{-n+\frac{2n\lg n}{w}
\lg\nparen{\frac{ew}{\lg n}}}.
$$
\end{lemma}

\subsection{Proof of the transition estimate}

Fourier inversion gives
$$
p_\ell=\E_y\hat\mu(y)^\ell.
$$
It follows from the identity in the preceding section that
$$
\Pr[W_\tau=e_{n+1}]
=e^{-\tau}\E_y\sinh\nparen{\tau\hat\mu(y)}.
$$
Moreover,
$$
\E_y\hat\mu(y)=\mu(0)=0.
$$
Set
$$
s(z)\defeq\sinh z-z.
$$
We therefore have
$$
\Pr[W_\tau=e_{n+1}]
=e^{-\tau}\E_y s\nparen{\tau\hat\mu(y)}.
$$

We shall use
$$
s(z)\le0\quad(z\le0),
\qquad
s(z)\le(\sinh1-1)z^3\quad(0\le z\le1),
\qquad
s(z)\le\frac12e^z\quad(z\ge1).
$$
Since $\abs{\hat\mu(y)}\le1$, it follows that
$$
\Pr[W_\tau=e_{n+1}]
\le
(\sinh1-1)e^{-\tau}\tau^3
\E_y\abs{\hat\mu(y)}^3
+
\frac12\Pr_y[\hat\mu(y)>1/\tau].
$$

Since
$$
\tau<R\le n\ln2<n,
$$
the first term is
$$
O\nparen{\frac{e^{-\tau}}V}.
$$
Since $\tau=\Theta(n)$, for all sufficiently large $n$ this is
less than
$$
\frac{\tau e^{-\tau}}{2V}.
$$

For the second term, put
$$
\theta\defeq
\frac{2\lg n}{w}\lg\nparen{\frac{ew}{\lg n}}.
$$
Our assumption $w\gg\lg n$ implies $\theta\longrightarrow0$.
The same lemma gives
$$
\Pr_y\sparen{\hat\mu(y)>1/\tau}
\le
\Pr_y\sparen{\abs{\hat\mu(y)}>1/n}
\le2^{-n(1-\theta)}.
$$
Since $R=\Theta(n)$, for all sufficiently large $n$ we have
$n\theta\ln2\le\zeta R/2$.  Therefore
$$
\begin{aligned}
\frac{\frac12 2^{-n(1-\theta)}}
     {\tau e^{-\tau}/V}
&=
\frac1{2\tau}
\exp\left(-R+\tau+n\theta\ln2\right)\\
&\le
\frac1{2\tau}e^{-\zeta R/2}
<\frac12
\end{aligned}
$$
for all sufficiently large $n$.  This proves
the desired transition bound, and hence $\E Z<1$.

Some choice of $M$ therefore has no odd circuit in $\kero$.
Appending the corresponding row $v$ described in the preceding
section produces a distance-$d$ code of size at least
$$
2^{n-m-1}
=\frac{\tau2^n}{2V}
\ge
\frac{(1-\zeta)R}{4}\frac{2^n}{V}.
$$
Thus, for every fixed $0<\zeta<1$ and all sufficiently large $n$,
there is a code $C$ satisfying
$$
\size{C}
\ge
\frac{1-\zeta}{4}
\ln\nparen{\frac{2^n}{V}}\frac{2^n}{V}.
$$
As $\zeta$ can be chosen arbitrarily small, there are codes for
which
$$
\liminf_{n\longrightarrow\infty}
\frac{4V\size{C}}{2^n\ln(2^n/V)}\ge1.
$$
For example, if $d/n\longrightarrow1/5$, then the multiplicative
improvement over the Gilbert--Varshamov bound satisfies
$$
\liminf_{n\longrightarrow\infty}
\frac1n\frac{\size{C}}{2^n/V}
\ge
\frac{\ln2}{4}\nparen{1-\entb{1/5}}
=0.048\ldots.
$$

\begin{remark}
\label{rem:linear-independence}
A slight modification also shows that $\kero$ is linearly
independent, with slightly larger numerical constants in the
estimates but no loss in the leading asymptotic constant.
Let $Z_{\mathrm{all}}$ count all circuits in $\kero$, and put
$\phi(z)=e^z-1-z-z^2/2$. The circuit count and Fourier
inversion give
$$
\E Z_{\mathrm{all}}
\le\frac V\tau\sum_{\ell\ge3}\frac{\tau^\ell}{\ell!}p_\ell
=\frac V\tau\E_y\phi\nparen{\tau\hat\mu(y)}.
$$
We have $\phi(z)\le0$ for $z\le0$,
$\phi(z)\le(e-5/2)z^3$ for $0\le z\le1$, and
$\phi(z)\le e^z$ for $z\ge1$. Splitting at
$\hat\mu(y)=1/\tau$ and applying
\autoref{lem:ball-fourier-estimates} as above yields
$$
\begin{aligned}
\E Z_{\mathrm{all}}
&\le(e-5/2)V\tau^2\E_y\abs{\hat\mu(y)}^3
  +\frac{Ve^\tau}{\tau}\Pr_y[\hat\mu(y)>1/\tau]\\
&\le O(1/\tau)+\frac{e^{-\zeta R/2}}{\tau}
=o(1).
\end{aligned}
$$
Thus $\kero$ is linearly independent with probability $1-o(1)$
for the same choice of $m$.
\end{remark}
